% Section 2.1, from lectures/L02/appendix/a_registers_and_io.md.

\section{Registers and I/O}
\label{c2:sec:registers}\label{c2:app:a}

\subsection{What's different about embedded}
\label{c2:a:1}
\begin{itemize}
  \item A microcontroller (MCU) is a complete computer with no screen or keyboard; I/O happens
    through pins wired to sensors and actuators (buttons, LEDs, motors, sensors).
  \item Your code is still compiled to machine code, but there's no OS between you and the
    hardware; you write directly to memory-mapped registers.
  \item C is the default choice here: predictable, tiny memory footprint, close to the metal, and
    still far more productive than assembly. A very skilled assembly programmer can occasionally
    beat a modern compiler, but rarely by much, and it's almost never worth it.
\end{itemize}

\subsection{The board: Arduino Uno/Nano}
\label{c2:a:2}
Either board works for this book: same \textbf{ATmega328P} chip, same registers, same
PIN$\leftrightarrow$PORT mapping for D0--D13 and A0--A5. The Nano just adds two extra analog-only
pins (A6, A7) and drops the barrel power jack.

The ATmega328P is an 8-bit AVR microcontroller with three I/O ports, each covering a group of pins:

\begin{narrowtable}{@{}ll@{}}
  \thead{Board pin} & \thead{Register port} \\
  \midrule
  D0--D7 & \code{PORTD0}--\code{PORTD7} \\
  D8--D13 & \code{PORTB0}--\code{PORTB5} \\
  A0--A5 & \code{PORTC0}--\code{PORTC5} \\
\end{narrowtable}

\bookimage[0.3]{lectures/L02/appendix/images/arduino_uno_pinout.png}{The Arduino Uno: digital pins
D0--D13 along the top edge, analog pins A0--A5 along the bottom.}{c2:fig:uno}

\noindent Here's the Nano for direct comparison: different board outline and a Mini-USB connector
instead of the Uno's USB-B, but the same D0--D13/A0--A5 pin numbering and the same ATmega328P
underneath:

\bookimage[0.45]{lectures/L02/appendix/images/arduino_nano_pinout.png}{The Arduino Nano: the same
pin numbering, plus the analog-only pins A6 and A7.}{c2:fig:nano}

\noindent\code{PORTB} and \code{PORTD} are purely digital I/O. \code{PORTC} is the analog-capable
port (used with the ADC, Chapter~\ref{c4:ch}) but behaves as ordinary digital I/O when you're not
doing analog reads.

For consistency, every breadboard circuit diagram in this section (and in \secref{c2:app:c}) shows
an Arduino Uno. If you're using a Nano, wire it the same way: same chip, same D0--D13/A0--A5 pin
numbering, the only difference is the board outline and USB connector in the picture.

\subsection{Bitwise operations (this is the whole game)}
\label{c2:a:3}
There's no \code{digitalWrite()} here; you write directly to registers, and a register is just a
byte, so every pin operation is a bitwise operation on that byte.

\begin{booktable}{@{}p{5.4em}p{6.4em}L@{}}
  \thead{Operation} & \thead{C operator} & \thead{Effect} \\
  \midrule
  OR & \code{|} & sets one or more bits, leaves the rest (``bitwise addition'') \\
  AND & \code{&} & reads a bit, or (with NOT) clears one (``bitwise multiplication'') \\
  NOT & \code{~} & inverts every bit \\
  XOR & \code{^} & clears matching bits, sets differing bits \\
  NOR & \code{~(a | b)} & inverse of OR: bit set only where both operands are 0 \\
  NAND & \code{~(a & b)} & inverse of AND: bit set unless both operands are 1 \\
  XNOR & \code{~(a ^ b)} & inverse of XOR: bit set where operands match \\
\end{booktable}

\noindent There's no single C operator for NOR/NAND/XNOR; you get them by combining \code{~} with
\code{|}/\code{&}/\code{^} as shown above. Don't confuse NAND with the clear-bit idiom below:
\code{reg &= ~mask} is AND with a \emph{negated mask}, not a NAND of \code{reg} and \code{mask};
those are different operations that happen to both involve \code{~} and \code{&}:
\begin{itemize}
  \item Set bit 3 without touching the rest: \code{reg |= 0b00001000;}
  \item Clear bit 3 without touching the rest: \code{reg &= ~0b00001000;}
\end{itemize}

\subsubsection{Shifting}
\code{1U << n} puts a single \code{1} in bit position \code{n}; this is how you address ``which
pin'' without hand-writing the bit pattern:

\begin{ccode}
#define LED1 1U         // D9 -> PORTB1.
PORTB |= (1U << LED1);  // Enable LED1.
PORTB &= ~(1U << LED1); // Disable LED1.
\end{ccode}

\noindent\code{x << n} and \code{x >> n} are also multiply/divide by \code{2^n}, if you ever need
that framing.

\subsection{Driving an output, reading an input}
\label{c2:a:4}
\begin{itemize}
  \item \code{DDRx} (\textbf{D}ata \textbf{D}irection \textbf{R}egister) decides, per bit, whether a
    pin is an output (\code{1}) or input (\code{0}).
  \item On an output pin, \code{PORTx} sets the pin high/low.
  \item On an input pin, \code{PORTx} instead enables (\code{1}) or disables (\code{0}) the internal
    pull-up resistor. With it enabled, the pin reads a clean \code{1} when unconnected/floating and
    \code{0} when pulled to ground, instead of floating unpredictably.
  \item To read the actual input level, read \code{PINx} (not \code{PORTx}).
\end{itemize}

\begin{ccode}
#define BTN1 5U                               // D13 -> PORTB5
#define BTN1_PRESSED (PINB & (1U << BTN1)) // High if BTN1 is pressed, low otherwise.
\end{ccode}

\noindent The button circuit used in \secref{c2:a:5} and \secref{c2:a:6} doesn't actually rely on
this internal pull-up: it wires an \textbf{external pull-down} resistor instead (one leg of the
button to 5\,V, the other leg to D13 and, through a 10\,k\textohm{} resistor, to GND). That means
D13 reads \code{0} when idle and \code{1} when pressed, the opposite of the internal-pull-up pattern
described above. The code still enables the internal pull-up anyway (\code{PORTB = (1U << BTN1);}),
so you see \code{PORTx}'s input-mode behavior used directly, but it isn't what makes this circuit
work. It's harmless in practice: the internal pull-up and the external 10\,k\textohm{} pull-down
form a voltage divider, and with the internal resistor's typical value of around 35\,k\textohm{}
the idle pin sits near 1.1\,V, comfortably below the 1.5\,V ($0.3 \times 5$\,V) under which a pin
is guaranteed to read \code{0}. It isn't guaranteed, though: the datasheet only promises
20--50\,k\textohm, and at 20\,k\textohm{} the idle pin rises to about 1.7\,V, where the reading is
undefined. In a design of your own, don't enable a pull-up against an external pull-down.

\subsection{First program: LED follows a button}
\label{c2:a:5}
An LED is connected to D9, and a button is connected to D13. The button is continuously polled:
\begin{itemize}
  \item The LED is enabled if the button is pressed.
  \item The LED is disabled if the button isn't pressed.
\end{itemize}

\bookimage{lectures/L02/appendix/images/circuit_led_follows_button.png}{An LED on D9 and a button on
D13, with its 10\,k\textohm{} pull-down resistor to GND.}{c2:fig:ledbutton}

\begin{ccode}
#include <avr/io.h>

/** GPIO pins. */
#define LED1 1U // D9  -> PORTB1.
#define BTN1 5U // D13 -> PORTB5.

/** GPIO operations. */
#define LED1_ON PORTB |= (1U << LED1)         // Enable LED1.
#define LED1_OFF PORTB &= ~(1U << LED1)       // Disable LED1.
#define BTN1_PRESSED (PINB & (1U << BTN1)) // High if BTN1 is pressed, low otherwise.

/**
 * @brief Set up system.
 */
static void setup(void)
{
    // Configure LED1 as output.
    DDRB = (1U << LED1);

    // Configure BTN1 as input with its pull-up enabled.
    PORTB = (1U << BTN1);
}

/**
 * @brief Application entry point.
 *
 * @return 0 on termination of the program (should never occur).
 */
int main(void)
{
    setup();

    while (1)
    {
        // Read BTN1, enable LED1 if pressed, disable if not.
        if (BTN1_PRESSED) { LED1_ON; }
        else { LED1_OFF; }
    }
    return 0;
}
\end{ccode}

\noindent\code{LED1_ON}, \code{LED1_OFF}, and \code{BTN1_PRESSED} are the object-like macros covered
in \secref{c1:a:12}: plain text substitution, with no parameters. \code{BTN1_PRESSED} is used as a
value, inside a condition, so its whole body is parenthesized and it behaves correctly no matter
what surrounds it at the call site. \code{LED1_ON} and \code{LED1_OFF} are complete statements,
only ever used on their own, followed by a \code{;}.

You'll see a Doxygen comment above every non-trivial function throughout this codebase, written in
the imperative mood (``Configure\dots'', ``Compute\dots'', not ``Configures\dots'' or ``This
function configures\dots'').

\subsection{Blinking two LEDs, with a hand-rolled delay}
\label{c2:a:6}
LEDs on D9/D10 alternate every 200\,ms while the button on D13 is held; both off otherwise.

\bookimage{lectures/L02/appendix/images/circuit_two_led_blink.png}{Two LEDs on D9 and D10, and the
button on D13.}{c2:fig:twoleds}

\begin{ccode}
#define F_CPU 16000000UL // CPU frequency in Hz.

#include <stdint.h>

#include <avr/io.h>
#include <util/delay.h>

/** GPIO pins. */
#define LED1 1U // D9  -> PORTB1.
#define LED2 2U // D10 -> PORTB2.
#define BTN1 5U // D13 -> PORTB5.

/** GPIO operations. */
#define LED1_ON PORTB |= (1U << LED1)      // Enable LED1.
#define LED1_OFF PORTB &= ~(1U << LED1)    // Disable LED1.
#define LED2_ON PORTB |= (1U << LED2)      // Enable LED2.
#define LED2_OFF PORTB &= ~(1U << LED2)    // Disable LED2.
#define BTN1_PRESSED (PINB & (1U << BTN1)) // High if BTN1 is pressed, low otherwise.

/** Time parameters. */
#define TICK_1MS 1U         // Generate 1 ms tick in delay_ms().
#define BLINK_SPEED_MS 200U // Blink speed in ms.

/**
 * @brief Set up system.
 */
static void setup(void)
{
    // Configure LED1 and LED2 as outputs.
    DDRB = (1U << LED1) | (1U << LED2);

    // Configure BTN1 as input with its internal pull-up enabled.
    PORTB = (1U << BTN1);
}

/**
 * @brief Generate delay.
 *
 * @param[in] ms Delay duration in ms.
 */
static void delay_ms(const uint16_t ms)
{
    for (uint16_t i = 0U; i < ms; ++i)
    {
        _delay_ms(TICK_1MS);
    }
}

/**
 * @brief Blink the LEDs.
 *
 * @param[in] blink_speed_ms Blink speed in ms.
 */
static void led_blink(const uint16_t blink_speed_ms)
{
    // Enable LED1, disable LED2.
    LED1_ON;
    LED2_OFF;
    delay_ms(blink_speed_ms);

    // Disable LED1, enable LED2.
    LED1_OFF;
    LED2_ON;
    delay_ms(blink_speed_ms);
}

/**
 * @brief Turn both LEDs off.
 */
static void leds_off(void)
{
    LED1_OFF;
    LED2_OFF;
}

/**
 * @brief Application entry point.
 *
 * @return 0 on termination of the program (should never occur).
 */
int main(void)
{
    setup();

    while (1)
    {
        // Read BTN1, blink the LEDs if pressed, disable both LEDs if not.
        if (BTN1_PRESSED) { led_blink(BLINK_SPEED_MS); }
        else { leds_off(); }
    }
    return 0;
}
\end{ccode}

\noindent A few things worth noting:
\begin{itemize}
  \item \code{F_CPU} must be \code{#define}d \emph{before} including \code{<util/delay.h>}, since
    the delay functions compute their loop counts from the clock frequency at compile time.
    \code{int} is only 16 bits on this target, far too small for 16{,}000{,}000; the \code{UL}
    suffix makes the constant an \code{unsigned long} explicitly, rather than leaving its type to
    C's rules for unsuffixed constants.
  \item \code{<stdint.h>} is included explicitly for \code{uint16_t}, even though \code{<avr/io.h>}
    happens to pull it in already; include what you directly use rather than relying on a header's
    own includes, which can change.
  \item \code{_delay_ms()} only accepts a compile-time constant or macro as its argument, not a
    variable, so it can't take \code{ms} directly. \code{delay_ms()} works around this by calling
    \code{_delay_ms(TICK_1MS)} in a loop, incrementing a local counter each pass until it reaches
    \code{ms}.
  \item \code{_delay_ms()} (and by extension \code{delay_ms()}) \textbf{blocks} the entire program
    while it waits, fine for this example, but this exact limitation is what the timers chapter,
    Chapter~\ref{c5:ch}, exists to solve.
\end{itemize}
